% ===================================================
% CHAPITRE 15 — BOURG, POPULATION, INDUSTRIES, ETC.
% Pages imprimées 75-95 — vues 78-98
% ===================================================
\chapter{Bourg, population, industries, etc.}

% --- Page imprimée 75 — vue 78
Bourg - Population
Divers quartiers
Industrie. Etc.

D'après le dernier recensement (10 mars 1946) la Commune d'Ouroux compte 509 habitants répartis entre 162 foyers. Le tableau ci-dessous montre qu'à la campagne quelques maisons sont inhabitées puisqu'il y a moins de ménages que de maisons. Il est à remarquer aussi que le nombre de ménages dans les divers quartiers de la Commune, tout en étant en nombre inférieur à ceux du bourg, fournit cependant une population plus forte que dans le bourg.

                              maisons ménages individus français étrangers.
Population éparse
(divers quartiers)              81      76       279       278        1
Population agglomérée
au chef-lieu (bourg)            71      86       230       228        2
Population totale              152     162       509       506        3

D'après le recensement de 1886 la Commune d'Ouroux compte 1000 habitants répartis entre 202 foyers. Le manuscrit de Lyon n'en mentionne en 1667 que 134 sans toutefois donner le chiffre de la population elle-même. Le cadastre de 1825 porte le nombre de maisons à 163. Mais depuis cette époque les grands propriétaires ont réuni à leurs fermes les diverses acquisitions qu'ils ont faites et un grand nombre de maisons ont été abandonnées puis démolies.

% --- Page imprimée 76 — vue 79
D'après une note conservée dans les registres de la paroisse la population était de 1018 âmes en 1831, de 978 en 1836, de 1098 dix ans plus tard. Le chiffre des naissances comme celui des décès est relativement élevé. De 1879 à 1889 la moyenne a été de 28,6 pour les naissances et de 21 pour les décès. Pendant cette même période la moyenne de la vie a été de 39 ans 11 mois.

Aucun document positif ne nous permet d'établir d'une manière certaine cette même population dans les siècles passés. Il est probable qu'elle n'a jamais dépassé le chiffre de l'année 1846. Au XVIIème siècle les 134 feux ne pouvaient guère fournir plus de 1000 habitants. Au besoin les rôles de grande taille suffiraient pour nous en convaincre. Celui de 1674 énumère 11 appentionnaires, 16 locataires, 40 grangiers ou fermiers, 3 tailleurs d'habits, 2 cordonniers, 1 maréchal, 1 mercier, 1 marchand et 1 charpentier. Ajoutez à ce nombre les exempts et nous ne trouverons pas un nombre de feux égal à celui de 1886.

Si la population de la campagne a diminué par suite de la fusion de la petite propriété dans la grande, celle du bourg s'est considérablement accrue au siècle précédent.

Le bourg d'Ouroux situé à 427 mètres d'altitude (seuil de l'église) occupe les deux rives de la Grosne bien que la partie principale soit adossée à la colline qui le domine au couchant. La rue la plus importante est formée par l'ancienne voie romaine dont nous avons parlé plus haut. Le bourg d'Ouroux qui compte actuellement 71 maisons, 86 foyers et 230 âmes, comptait il y a 70 ans : 62 maisons, 92 foyers et 375 âmes. Il comprend plusieurs petites industries qui attirent les populations voisines et lui donnent une aisance relative.

D'après une tradition rapportée par Mr Morétain, le bourg d'Ouroux à une époque reculée ne renfermait que six foyers. Nous croyons qu'il suffit de remonter aux premières années du XVIIIème siècle pour constater qu'il avait peu d'importance. Le plan du fief de Montolieu qui s'arrête à la voie romaine ne mentionne que deux maisons dans la grande rue actuelle : celle de Mr Canard et celle de MM Gonnet. Toutes les autres constructions sont de date récente. Nous n'aurions eu dès lors que la Cure, la maison Duligier (maison Joly,
% --- Page imprimée 77 — vue 80 (suite)
puis maison Rollet) et Bouchacourt, les deux citées plus haut et la maison Martin (aujourd'hui Aufranc, menuisier) près de la rivière. Une autre maison très ancienne était celle qui précédait la maison Desperrier, de date récente et qui forme l'angle entre la grande rue et la route de Fontmartin.

Dans son état actuel le bourg n'offre aucune construction qui mérite d'être signalée. Toutefois, à la suite du bombardement durant cette dernière guerre (dont nous parlerons plus loin) plusieurs maisons ont dû être reconstruites ou réparées et grâce à l'aide apportée par la "Reconstruction" la partie du bourg qui se trouve à l'intersection de la grande rue, de la route de Fontmartin et de celle de la Serve a été enjolivée.

La cure est vaste et doit être très ancienne. Aucun titre ne fait allusion aux constructions successives qui la composent. Si un monastère a existé jadis dans le jardin, nous serions portés à croire qu'on utilisa, après sa destruction, les bâtiments qui en dépendaient pour loger les vicaires perpétuels qui succédèrent aux religieux. Mr Bouteille, à son arrivée à Ouroux la trouva dans un état de délabrement impossible à décrire. Des réparations sagement conduites par lui et ses successeurs en ont fait un presbytère commode et agréable. Mr Morétain avait lui-même pris l'initiative en faisant construire la tour qui fait communiquer les étages avec le rez-de-chaussée et qui donne à la cure un certain cachet d'antiquité.

\begin{center}\textit{[Illustration : la cure et sa tour (vue prise du jardin)]}\end{center}

À signaler encore la maison Gonnet, touchant la tour de Nagu, qui a dû être construite avec des matériaux ayant appartenu primitivement au château de Nagu. La fenêtre du premier étage (dans la cour intérieure) offre tous les caractères de la Renaissance : linteau en accolade, gorge profonde à l'angle extérieur des pieds droits dans laquelle sont sculptées des tulipes. Des pierres semblables et appartenant à d'autres fenêtres ont été trouvées en
% --- Page imprimée 78 — vue 81 (suite)
fouies et sont maintenant conservées dans la même maison. D'après
le témoignage des anciens, ces pierres sculptées faisaient partie de
fenêtres avec meneau; les débris qui subsistent et qui sont ornés
sur les deux faces viennent confirmer la tradition.

Nous signalerons en passant une particularité curieuse qui se
trouvait dans les anciennes maisons de la campagne. Une seule, construite
dans ces conditions se voyait encore à Cenves il y a une cinquantaine
d'années. Le foyer était au milieu même de la cuisine
dont le sol était formé par de la terre battue. Quatre fortes poutres
de chêne, formant un carré, supportaient deux paravents que l'on
soulevait, à volonté, de l'intérieur et qui formaient ainsi une cheminée
faisant saillie au-dessus de la toiture. Aujourd'hui, la plupart
des fermes sont construites selon toutes les formes architecturales.

Quel fut, dans les siècles passés l'importance de l'industrie du
bourg d'Ouroux?

Examinons d'abord la valeur historique de cet ancien dicton journellement
répété: "ville d'Ouroux, bourg de Beaujeu". Un vieux titre
sans date, au rapport de Mr de la Carelle, conservé longtemps aux archives
de Villefranche et présentant les caractères d'un manuscrit
du XIème siècle, donnait à Ouroux le titre de "ville". Notre historien
ignorait certainement les Chartes qui nous ont été conservées
par le Cartulaire de St-Vincent de Mâcon et dans lesquelles Ouroux
est également désigné sous le nom de "villa oratorii"

Le mot latin "villa" doit se traduire par village et non par
ville. La ville n'était pas simplement une agglomération d'habitations;
elle comprenait aussi le territoire entourant le village.
D'après Du Cange, le mot "burgus" désignait un village construit sous
les murs d'un château féodal. Or Beaujeu était dans ce cas puisque
les Seigneurs de Pierre Aiguë apparaissent dès la fin du Xème siècle.

Comme on le voit, une fausse interprétation du mot "villa" a déterminé
les historiens à attribuer à notre village une importance
qu'il n'a jamais eue.

% --- Page imprimée 79 — vue 82
Devons-nous citer ici l'opinion de Mr Morétain relative à l'industrie de notre vallée ? "À une époque, dit-il, qui ne doit pas remonter au-delà du XIIIème siècle, il se fabriquait à Ouroux, je ne sais quelle marchandise en métal qui se poinçonnait. On a retrouvé trois de ces poinçons formant un ovale rempli de quelques figures et de lettres ; l'un d'eux portait les armoiries d'une famille inconnue. Peut-être serait-ce les matières que l'on fabriquait à la place du lieu qu'occupe le presbytère où il paraît avoir eu des fourneaux à métaux, chauffés avec du bois, car en creusant des fondations, il y a quelques années, nous avons trouvé une immense quantité
\begin{center}\textit{[Illustration : Coin Sud-Est. - Grande rue d'Ouroux et arrivée du car Tramayes-Belleville-Lyon]}\end{center}
de cendres mêlées à des scories de métaux englobant des charbons de bois. On voit des pierres fondues, vitrifiées, tordues par l'ardeur du feu. Les poinçons dont j'ai parlé servaient peut-être à marquer ces métaux". Nous n'avons pu malheureusement retrouver ces poinçons dont nous parle Mr Morétain. Leur étude eût servi certainement à contrôler son assertion.

Le même auteur raconte qu'à la même époque, en déracinant un vieil arbre dans le pré Duligier, au-dessus de la maison Commune actuelle, on trouva dans une enceinte circulaire qui paraissait être une ancienne tour, cinq fours dont l'un était fort grand et les autres avaient plus ou moins d'étendue. Non loin de là on découvrit une forge presque entière avec ses soufflets et autres ustensiles.

% --- Page imprimée 80 — vue 83
Au temps des Seigneurs de Nagu le village comprenait les quelques maisons construites au pied du château. Mais lorsque le fief fut démembré, les divers acquéreurs qui appartenaient à des familles aisées vinrent s'établir au centre de leurs acquisitions. Nous avons déjà nommé les De Bosco, les De Gorze. Leur présence suffit à donner un nouvel élan à la population et bientôt les échanges entre les localités voisines se firent à Ouroux. Ainsi par lettres patentes du 1er Octobre 1523 données par François 1er et confirmées par le roi Charles IX, un marché fut établi le lundi de chaque semaine ainsi que deux foires dans l'année, celle de la Saint-Antoine et celle de la St-Barthélemy. Seule la première existe encore actuellement. Trois autres ont été établies plus tard : le 29 Avril, le 20 Juin et le 7 Novembre. Ces trois dernières foires n'avaient aucune importance vers la fin du siècle dernier et actuellement n'existent plus. Le marché du lundi n'existe plus depuis longtemps. On n'a même conservé aucun souvenir de son établissement.

Il est probable que la construction de la digue élevée au lieu dit : « La Croix de l'Étang, » remonte à cette époque. Nous ne la trouvons mentionnée qu'une seule fois dans les actes en 1623, lorsque Philothée de Lafont donne à son frère, curé de Verfe, une maison située au bourg d'Ouroux sur la grande rue tenant du côté des Chassignes à la Croix du Rempart d'Ouroux. C'est là du moins la seule interprétation que nous puissions donner à ce texte.

Nous avions cru voir également cet étang mentionné dans la copie du terrier d'Illoisnet (cote 79 vo) à la reconnaissance de Guy Jocond dans laquelle on parle de l'escluze. La reconnaissance précédente de Pierre Jocond nous montre qu'il s'agit de l'escluze d'un moulin situé sur la rivière du Juel au point de jonction du chemin de Jorlet avec celui de Braille.

La tradition rapporte qu'un membre de la famille de Gorze, celui qui se rendit acquéreur d'une partie du fief de Nagu, vint à Ouroux comme tanneur. Le pré acheté par la fabrique et situé sur la rive droite de la Grosne, près de la cure, a conservé du reste le nom de tannerie. À ce moment même on établit un marché à Ouroux, on crée
% --- Page imprimée 81 — vue 84 (suite)
deux foires, chose rare à cette époque. Notre bourg devient ainsi un centre industriel et commercial. Mais, pour alimenter les usines qui viennent d'être construites, il faut une eau abondante et, nous l'avons vu, la Grosne est souvent à sec en été. Comment parer à cette difficulté ? Nos ancêtres mirent à exécution un projet fréquemment réalisé de nos jours. A 1 kilomètre et demi du bourg la vallée est très resserrée. Rien de plus facile que d'établir un barrage qui formerait ainsi un vaste réservoir dont on pourrait utiliser l'eau dans les moments de sécheresse.

On construisit donc un mur très solide de 1 m 30 d'épaisseur, allant d'une montagne à l'autre. On amena une quantité considérable de terre, prise sur la montagne qui domine la digue du côté du midi et qui devait servir de soutènement et de contrefort à ce mur. Au centre même de la digue on plaça, au niveau du sol une "chasné" en chêne qui n'a disparu que dans ces dernières années. L'ancienne voie étant désormais submergée, un sentier fut tracé sur le flanc de la montagne. Son point de départ était près du chemin qui va de la route à Braille et il est facile d'en suivre le tracé à travers les rochers, tracé qui nous permet de fixer approximativement l'étendue de l'étang.

Qu'arrive-t-il ? Nos ingénieurs avaient compté sans l'énorme poussée produite par cette masse d'eau considérable. La digue ne put résister et creva vers le milieu. Six mois à peine, un an peut-être dut suffire pour que l'eau, s'infiltrant, détrempât le terrain. Le mur lui-même était construit dans des conditions trop défavorables pour qu'il pût résister. Les blocs énormes de pierres qui existent encore et servent de soubassement à la nouvelle route, prouvent que les constructeurs avaient fait leur possible pour parer à pareille éventualité. L'expérience leur manquait encore !

Nous pouvons juger de l'état de la digue par les restes qui ont été conservés. Le mur, quoique formé de pierres légères noyées dans le mortier est très solide. Le terrain qui n'a pas été entraîné, par les eaux, forme comme une espèce de tertre qui est cultivé, comme la montagne elle-même. Mais peu à peu la muraille disparaît, sous l'influence du temps et avec elle, le tertre dénué par l'éco
% --- Page imprimée 82 — vue 85 (suite)
buage, voit, chaque année sa hauteur diminuer sensiblement. Bientôt la "digue de l'étang" ne sera plus qu'un souvenir.

\begin{center}\textit{[Illustration : "la croix de l'étang"]}\end{center}

Il est facile de comprendre quel dut être le sort de la partie basse du bourg d'Ouroux lorsque les eaux se précipitèrent par la trouée qu'elles venaient de se former. On doit évidemment rapporter à cette époque la date du désastre dont on a gardé le souvenir et dans lequel le bourg fut détruit par les eaux. Mr Morétain cite le fait d'un meuble rempli de linge et qu'on a retrouvé enfoui profondément. L'imagination populaire s'est donné beau jeu au sujet du cataclysme et parfois même on a fait intervenir l'action diabolique. Le site sauvage de "la croix de l'étang" favorisait merveilleusement la verve des narrateurs. Mr Bouteille a fait placer sur le rocher qui domine les restes de la digue la croix que Mr Trouilloux avait fait élever au sommet du clocher après la Grande Révolution.

Ouroux, nous dit encore la tradition, fut dans les siècles passés détruit par le feu. Le désastre dut, sans doute, avoir lieu à l'époque de la destruction du château de Nagu par les Tard-Venus dont nous avons parlé précédemment.

Des incendies considérables sont venus récemment modifier la situation générale du bourg. Nous signalerons plus particulièrement ceux de 1867, 1870 et 1873. Dans le premier, la maison du bureau de bienfaisance occupée alors par les religieuses fut sur le point de devenir la proie des flammes ; les maisons voisines furent détruites.

C'est ici le lieu de rappeler l'ouragan qui eut lieu le 20 Février 1879 et causa à Ouroux des dégâts incalculables. Un cyclone, d'une violence inouïe, qui se forma dans le centre des Cévennes, pour aller se perdre dans les montagnes de la Suisse, traversa
% --- Page imprimée 83 — vue 86 (suite)
notre région en s'attaquant de préférence aux arbres les plus beaux et les plus anciens. Le bois du Thel, ceux de l'Hospice, de Lacarelle du bureau de bienfaisance furent gravement endommagés. Toute la partie nord de la forêt de Gros-Bois fut renversée et du bois légué par Mr Guillin à la fabrique il ne subsistait que trois sapins. Un quart d'heure à peine suffit pour tout détruire malgré la résistance que devaient offrir des sapins ayant près d'un siècle.

Ouroux, par sa situation sur la voie reliant Paris au midi de la France devait payer son tribut à toutes les grandes invasions, comme à tous les fléaux qui en sont les conséquences.

Nous savons que les Protestants s'étant rendus maîtres de Mâcon, de Cluny et de tous les châteaux-forts de la vallée, vinrent à Ouroux, profanèrent l'église, scène qu'un curé de la paroisse avait fait représenter par la peinture sur la façade de l'église. C'est le seul souvenir qui nous ait été conservé de leur passage parmi nous.

Mais notre population eut à souffrir surtout pendant les guerres de Ligue. En 1589 Beaujeu avait embrassé le parti des ligueurs.

François de Nagu de Varennes fut nommé gouverneur du château. La soldatesque effrénée qui composait la garnison parcourait le pays, enlevait les récoltes, dévalisait les maisons et avait, chaque jour, avec la population, des démêlés qui finissaient toujours par le meurtre et l'incendie. Cet état de choses dura jusqu'en 1594, époque où Lyon se soumit à Henri IV.

Ajoutez à cela le passage incessant des troupes plus ou moins indisciplinées, sous la conduite de commissaires, logeant chez l'habitant, se procurant elles-mêmes les vivres dans les localités qu'elles traversaient. On comprendra dans quelle situation devait se trouver Ouroux. Le rendement de Comptes fait par Antoinette De La Font, mère de Pierre Testenoire nous offre à ce sujet de curieux détails.

Après avoir énuméré divers objets comprenant le mobilier et en particulier des hallebardes, des masses d'armes, etc... elle ajoute :
"lesquels objets ont été prins, désrobés, emportés, rompus, brisés et bruslés par les gens de guerre qui ont logé par ces derniers troubles au dict lieu d'Oroux, mesme ceulx du Capitaine Bataille qui logèrent au dict Oroux en l'an 1590... ceux du Capitaine La Peyrouse,"
% --- Page imprimée 84 — vue 87 (suite)
qui y furent logés en l'an 1591... ceux du Capitaine La Rivière qui y furent logés en l'an 1594... ceux de l'armée du Roy conduits par le Sieur maréchal de Biron dict la Carrabine en l'année suivante.. ceux de Mr de Guise y furent logés la même année, environ la Saint-Michel. En l'année 1596 ceux du régiment du Sr de Nofay.. En l'année suivante y furent logés d'autres qui emportèrent..."

On peut se faire une idée de l'état précaire dans lequel devaient se trouver les habitants en des démarches qu'ils faisaient, pour éloigner ces bandes indisciplinées. Encore devaient-ils acheter leur départ par des dons ou des contributions volontaires. C'est ce que nous pouvons constater en étudiant les dépenses faites au nom de Pierre Testenoire. En 1587 le Sr de Nagu de Varennes ayant détourné une compagnie de gens de guerre qui voulaient loger à Ouroux, "les habitants du dict lieu, pour récompense, lui auraient fait un don de certaine somme de deniers qui furent levés sur les plus aysés de la dicte paroisse et en paya le dict Testenoire vingt trois sols". En l'année 1589, même récompense accordée au Sieur Despré ; la part du jeune Testenoire n'est que de 10 sols. En 1590 sa part est de 25 sols pour le Sr de Varennes. L'année suivante, le Sr Despré reçoit deux nouvelles gratifications. Pierre Testenoire est taxé, la première fois à 20 sols, la deuxième à 20 sols.

Nos familles "aysés" se lassèrent enfin et les déprédations suivirent leur cours.

Depuis cette époque notre bourg semble avoir moins souffert des guerres et des invasions qui ont eu lieu. En 1814 un détachement de corps d'armée Autrichien qui venait de s'emparer de Lyon, traversa nos montagnes et s'établit pendant une quinzaine de jours à Ouroux, mais les biens et les propriétés furent respectés.

Il n'en fut pas tout-à-fait de même au château Guillin à Avenas. C'est ce que nous apprend une lettre du curé d'Avenas adressée à Madame de Guillin d'Avenas habitant 2 Place St-Jean à Lyon

"Avenas 2 Avril 1814

"J'ai reçu votre lettre Jeudi. Je suis descendu au château où s'est trouvé Claire Champagnon. Je vous dirai donc Madame que le 10 Mars il arriva vers les 4 heures du soir environ 1200 hommes tant à
% --- Page imprimée 85 — vue 88 (suite)
pied qu'à cheval. Un assez grand nombre alla à Beaujeu le même soir et aussi aux Ardillats. Ces 1200 hommes voulurent loger à Avenas. C'était un régiment de dragons qui nous embarrassait le plus. Le chef s'adressa à moi et me demanda s'il lui serait possible de loger à Avenas tous ses hommes. Je lui dis que la Paroisse était trop petite et que c'était impossible. Il me répondit qu'il y avait un château où il en pourrait loger une partie. Je lui répondis qu'il était fermé, qu'il n'y avait personne. Alors il me demande de le faire conduire et selon toute apparence il se fit conduire au château avec au moins 180 hommes et autant de chevaux. Le foin qui était dans votre fenil ne fut pas épargné. Les écuries, toute la cour étaient pleines de chevaux et de foin. Il n'y a pas de portes qu'on ait ouvert : cave, chambres, armoires, tout fut visité ; tout le vin a été bu ; les serrures des portes et celles des armoires arrachées, les fusils emportés ; café, sucre, eau de vie, argenterie, linge furent emportés. On a fait aussi des réquisitions de vaches. Votre domestique n'a presque rien sauvé. On ne lui a pas laissé une chemise. À la Cure j'ai logé de 12 à 13 officiers dont j'ai été content. Ils m'ont bu assez de vin et à leur départ qui fut le 11 à 3 heures du soir ils me laissèrent une grande jument qui était trop fatiguée pour s'en servir."
A. Collonge, curé d'Avenas.

\begin{center}\textit{[Illustration : Armes de Guiffin de Pougelon seigneur d'Avenas De gueules, à 4 flèches d'argent formant le giron.]}\end{center}

La guerre de 1870 fut désastreuse pour notre paroisse. Les Mobiles furent plus particulièrement éprouvés dans les combats de Nuits où plusieurs trouvèrent la mort. Ceux qui ont échappé aux balles ennemies ont eu à endurer les privations et les souffrances les plus grandes durant la captivité en Allemagne.

Avant de terminer le sujet qui nous occupe, nous ne saurions passer sous silence un document intéressant et qui nous indique dans quelles conditions se recrutaient les milices pendant les guerres de Louis XIV. Les hommes non mariés étaient invités, au prône de la paroisse, à se réunir en assemblée le Dimanche suivant, à l'issue des Vêpres, afin de tirer au sort pour déterminer
% --- Page imprimée 86 — vue 89 (suite)
quel serait celui qui servirait dans les troupes de sa Majesté pendant 5 ans. Le notaire était présent et rédigeait l'acte, séance tenante.

En 1701 Guillaume Gobet s'offrait à partir moyennant une certaine somme d'argent qui fut débattue et fixée à 100 livres et, qui devait lui être payée lorsqu'il avait reçu l'ordre de se mettre en marche. Antoine Large, principal leveur des tailles devait en outre lui donner jusqu'au jour de son départ la somme de 4 sols par jour pour sa subsistance

Il ne sera pas sans intérêt d'étudier brièvement les transformations successives des noms donnés aux divers quartiers de notre paroisse. Aujourd'hui, encore, bien que fixés par les Actes civils et le cadastre, ils tendent à se modifier peu à peu par l'usage.

Sauf de rares exceptions, les fermes ont pris le nom du propriétaire ou des fermiers qui les ont occupées. Nous l'avons vu pour La Carelle que l'on désignait communément par "La Combe des Essers ou des Essarts". La vallée prenait ensuite le nom du ruisseau appelé "les Loittes". L'orme de Crie, Cuisset, la Pierre du Sée, Fougères se trouvent mentionnés dans les titres du XVe siècle.

Le village de Chambuerd était connu à cette époque sous le nom de "Campo benedicto" (champ béni) ou peut-être champ à Benoit.

Les Pourchoux, autrefois La Condamine, appartenaient au commencement du XVIe siècle à Guichard au Porcheur, surnom que la propriété a gardé jusqu'à présent.

Laye, Le Planet, le Thel figurent dans les titres les plus anciens. Au XVIIe siècle, les Monnet portaient le nom de Mas Bleton et les Rochettes celui de Mas de Vaux. La Combe actuelle avait pris le nom de son propriétaire et s'appelait "La Combe Iouand" comme Jorlet "La Combe Genevret". Une famille Denis vint s'établir à la Combe et donna son nom à l'une des fermes.

Les Friats, Vigue, Dombay, les Mélinands remontent au XIVe ou au XVe siècle. La vallée des Jambons s'appelait jadis "Mas Dubost". La ferme Berthu du nom de Berthoud fermier au XVIIIe siècle fut
% --- Page imprimée 87 — vue 90 (suite)
connue sous le nom de "chasteau en Montaigne" et "château Chevrier" du nom de son possesseur.

Si nous remontons au delà du XVème siècle, la formation des noms obéit à des lois complètement différentes; les fermes ne prennent pas le nom de leurs possesseurs, mais ceux-ci adoptent celui des terres ou des fonds qu'ils exploitent. Primitivement l'individu n'avait qu'un nom, celui qu'il avait reçu au Baptême et que l'on faisait suivre du nom de son père, de son fils ou même de son frère. Plus tard les relations devenant plus fréquentes, ces désignations, vagues, ne pouvaient suffire. On ajouta au nom un surnom donné à l'occasion d'une habitude, d'un défaut, d'une difformité, du lieu de naissance. Le Cartulaire de Mâcon nous a conservé le nom de Joannes de Oratorio (Jean d'Ouroux). Plusieurs familles, dans notre contrée portent encore ce nom de Ouroux. Ceux de De Bosco, Du Bois ou Dubuis, Du Ligier, De La Font, Du Thel, Du Croux ou Du Creux n'ont pas eu d'autre origine. Malgré les transformations qu'ils ont dû subir il est facile d'en reconnaître l'origine.

\begin{center}\textit{[Illustration : un ouvrier paysan et paysanne à la fin du XVème siècle d'après des gravures du temps.]}\end{center}

% --- Page imprimée 88 — vue 91
Aux siècles derniers la division de la propriété différait
peu à peu de celle qui existe à cette heure. Nous tenons à conser-
ver le nom des divers propriétaires qui se partageaient la parois-
se vers 1650. Nous indiquerons l'état de leurs revenus à cette épo-
que.

Le Seigneur Comte de St-Maurice possédait au Mas Dubost (vallée
des Jambons) 5 domaines et un pensionnaire, le tout amodié au
prix de 1202 livres.

Le Seigneur de la Barmondière : 2 domaines, également au Mas Du-
bost et la rente noble d'Arcis : 550 livres

Le Seigneur d'Avenas : 3 domaines, 1 moulin et la moitié de son
grand pré de réserve du revenu de 900 livres. Il payait à la rente
noble d'Alloignet en 1650 :
        Argent........8 livres
        Froment.......24 quarterons, 3 coupons et demi
        Seigle........31       =       2       =      et demi
        Avoine........202      =       9       =      et quart
        Poulet........2 et un douzième d'autre
        Conil (lapin).. un demi

Le Seigneur de Bacot : 1 domaine et 1 moulin (actuellement mou -
lin Voland) le tout amodié au prix de 400 livres. Il payait à la
même rente :
        Argent........31 sols 2 deniers viennois
        Seigle........6 quarterons 8 coupons
        Orge..........2 raz
        Avoine........14 raz 1 quarteron, 3 coupons et 3/4
        Geline (poule). une demie

Le Seigneur de Gorze : 3 domaines dont un au Thel et un autre près
du bourg, une maison dans le bourg, 1 moulin (anciennement moulin Châ-
telet) avec les rentes nobles de Nay, Nagu et Sacconay qu'il avait
dans la dicte paroisse, le tout d'un revenu de 1200 livres.

Le Sieur de Verpré : son château de la Carelle, un pré de réserve,
7 domaines à la Combe des Essarts et aux Pourchoux du revenu de 1000
livres. En 1788 ses servis étaient de :
        Argent........8 livres, 13 sols, 4 deniers viennois
        Froment.......10 quarterons 6 coupons
        Seigle........38 quarterons
        Avoine........3 raz, 16 coupons
        Poule.........4 et une demie

% --- Page imprimée 89 — vue 92
Le Seigneur de Fautrières : 3 domaines amodiés au prix de 600
livres. En 1681 Madame de Fautrières payait à la rente d'Aloignet :
Argent.........5 livres 13 sols, 8 deniers
Froment........5 mesures, 9 coupons
Seigle.........29 mesures 6 coupons deux tiers
Avoine.........104 mesures
Géline.........1 et trois quarts d'autre

Le Sieur lieutenant d'Aigueperse bourgeois à Lyon : 3 domaines
(aujourd'hui les fermes de l'Hospice) du revenu de 600 livres. Il
payait : Argent.......5 livres 18 sols 2 deniers
         Froment.......1 quarteron 5 coupons
         Seigle........16    =     7    =
         Avoine........46  raz     4    =
         Géline........2 et demi

Laurent De La Font, officier de sa Majesté : 2 domaines dont l'un
au Carruge et l'autre à la Carelle (domaine Mestrat, aujourd'hui Allée
amodiés au prix de 525 livres

Le Seigneur de St-Julien (St-Mamert) : 1 domaine, 1 maison dans le
bourg du revenu de 150 livres. Il payait en 1788 :
         Argent viennois...43 livres 3 sols
         Froment...........12 coupons et demi, mesure de Beaujeu
         Seigle............4 quarterons, 5 coupons et demi
         Avoine............10 quarterons, 4 coupons et demi

Il devait encore à cause de ses fonds montant de la rente noble
d'Arcis en 1710 :
         Argent.......7 livres, 3 deniers
         Froment......2 tiers d'une coupe et 1 coupon
         Seigle.......10 coupons
         Avoine.......2 mesures 7 coupons

Nous passons sous silence les noms de ceux qui ne possédaient
qu'une maison et dont les revenus étaient peu importants.

Il est difficile de donner des détails précis sur la situation
financière de notre paroisse aux siècles passés. Les seuls titres
qui nous aient été conservés concernent l'époque où la famille Duligier
remplissait les principales charges de notre région. Nous
nous contenterons de signaler les quelques notes éparses que nous
avons pu recueillir à ce sujet.

Nous verrons plus bas que sous Guichard III (1200-1226) seigneur
de Beaujeu, la dîme d'Ouroux était évaluée à la somme de 100 sols.
Quel était le mode de perception employé parmi nous ? Il est probable
qu'il différait peu de celui qu'on suivait à Trades en 1769.

% --- Page imprimée 90 — vue 93
Nous citerons textuellement la déposition faite par les principaux habitants : "La dixme se prélève, savoir : de la quatorze la quinzième gerbe en seigle et de la quinze la seizième en froment, que la mesure se prélève avant la dite dixme et que lorsque cette prélévation n'est pas faite, l'on compte deux fois douze ou deux fois 13, ce qui fait que dans le nombre de gerbes de l'une ou l'autre qualité il s'en trouve une qui ne se dixme point".

En ce qui concerne Ouroux, nous n'avons pu relever qu'une seule note à ce sujet dans les papiers de la famille Duligier. Vers la fin du siècle dernier il prélevait pour la dîme : 24 gerbes de seigle sur son domaine de Chanteloup, 11 sur celui de Jorlet et 38 sur celui des Friats. Ces deux derniers domaines existent encore actuellement.

\begin{center}\textit{[Illustration : Beaujeu - Église Saint Nicolas (XIIe siècle)]}\end{center}

On peut se faire une idée approximative des charges qui pesaient sur les propriétés pour subvenir au service religieux.

Les principaux habitants affermaient du Chapitre de Beaujeu les dîmes de notre paroisse ; en 1685 une minute du rôle nous en énumère jusqu'à 19.

Quant au mode d'impôts ordinaires, qu'il nous suffise d'analyser un mandement de Louis XIV pour la perception des tailles, en 1693. Il nous fera connaître un des côtés et non le moins intéressant de la vie civile à cette époque.

"Les Consuls, disait l'ordonnance, doivent être élus le 1er Septembre. Seuls les Consuls, assesseurs et collecteurs ainsi que le greffier pouvaient assister à la confection des rôles et défense était faite à toute personne constituée en dignité, seigneur, gentilhomme, juge et officier de prendre connaissance de la confection des tailles, ni de s'immiscer ou d'empêcher les assemblées, sous peine de 500 livres d'amende.

Il était recommandé aux collecteurs de s'imposer, eux et leurs parents, comme tous les autres et de ne pas user de représailles con
% --- Page imprimée 91 — vue 94 (suite)
tre les collecteurs de l'année précédente ; leur cote ne devait être
augmentée qu'à proportion de l'augmentation des tailles.
    "Les ecclésiastiques étaient cotés pour les biens par eux acquis
purement et simplement ou reçus par succession collatérale ou dona-
tion. Ils ne l'étaient pas pour leurs biens d'église, leur titre pres-
bytéral, ni pour les successions en ligne directe. Leurs fermiers ne
devaient pas être exemptés, mais les prêtres pouvaient faire valoir
une ferme par eux-mêmes ou par des domestiques, sans être sujets à
la taille.
    "Les gentilshommes pouvaient pareillement tenir une ferme par
eux-mêmes ou par des domestiques sans être inscrits au rôle ; pour
les autres les fermiers étaient imposés d'après leur revenu.
    "Les collecteurs devaient confectionner les rôles et n'exempter
personne sous peine de 50 livres d'amende.
    "Les mineurs étaient exempts de la taille jusqu'à 18 ans.
    "Les habitants qui voulaient déloger d'une paroisse pour aller
dans une autre devaient faire publier à la Messe paroissiale et le
signifier aux habitants, procureurs, syndics, puis le faire enregistrer
au greffe de l'élection avant le 1er Octobre.
    "Les receveurs et collecteurs ne pouvaient faire saisir les ob-
jets suivants : bestiaux arables, pain, lit, habits et ustensiles ser-
vant à la culture des prés, vignes, terres ; ni les cuves, tonneaux,
pressoirs, tuiles, couverts de maison, portes, fenêtres ni outils des
artisans, ni la paille et le foin nécessaires pour la nourriture du
bétail arable, savoir : 8 charretés de 10 quintaux pour chaque paire
de bœufs et 6 chars pour chaque paire de vaches, outre le revivre.
Ils devaient encore laisser suffisamment de grain pour ensemencer
les terres.
    "Il était enjoint à toute personne de mettre en bon état les
chemins "chacun endroit soy" ; de même pour les ponts et passages,
sur les petites rivières, sous peine de 50 livres d'amende.
    En 1656 la généralité de Lyon fut imposée de 400169 livres 4 sols
6 deniers pour tailles, subsistance, droits d'officier, quartier d'hi-
ver, taillon et solde, ponts et chaussées, 12 deniers de l'armement na-
val, 6 deniers des directeurs de tailles, seconde partie de l'épar-
gne, réparation du palais et prison royaux de Lyon. Ouroux était im
% --- Page imprimée 92 — vue 95 (suite)
posé pour la somme de 2334 livres payables en deux termes. En 1788, ses impositions pour le même objet étaient de 2813 livres.

Il était recommandé aux collecteurs, lisions-nous plus haut, de ne pas user de représailles ; malheureusement ce conseil était fréquemment oublié. Nous en avons un exemple intéressant dans la famille Duligier. Quelques années après l'ordonnance royale que nous venons de rapporter, l'un des membres de cette famille adressait à la Sénéchaussée de Villefranche une requête à l'effet de faire diminuer ses tailles. Voici les motifs qu'il invoquait :

"Il dit qu'il a peu de biens dans les paroisses d'Ouroux et de Charentay, que la grêle et l'inondation leur ont fait beaucoup de mal en 1707. Ses biens fournissent à peine un revenu de 300 livres. Cependant il payait en 1703 52 livres de grande taille avec subsidiaux à proportion ce qui excédait de beaucoup la moitié de ses fonds et héritage.

"Cette cote excessive détermina Jean Chrysostome à se faire pourvoir de la charge de substitut du procureur du Roi en l'élection de Beaujolais, à cause du privilège accordé à cette charge de l'exemption des tailles et autres impositions. Pour effectuer le payement de cette charge, il emprunta 2500 livres. Le Roi ayant retiré le privilège de l'exemption, les consuls et habitants de la paroisse d'Ouroux portèrent ses tailles à 70 livres en 1705, à 100 livres en 1706 bien qu'il n'eût fait aucune acquisition nouvelle depuis 18 ou 20 ans.

En 1707, ayant appris que les Duligier voulaient se pourvoir contre ces diverses augmentations, les consuls, manants et habitants firent cotiser d'office Jean Hugues à 110 livres et Jean Chrysostôme à 55 livres bien que celui-ci eût reçu en donation tous les biens de son père. En 1708 ils furent imposés de 150 livres. Impossible de recueillir une pareille somme dans leurs fonds et ils demandaient à voir leurs tailles réduites de 50 livres.

"Le 6 Janvier 1708 après la messe et le peuple assemblé au son de la cloche en la manière accoutumée, devant la grande porte de l'église, comparurent Jean Hugues et Jean Chrysostôme Duligier en présence des Consuls et des collecteurs de tailles de l'année présente
% --- Page imprimée 93 — vue 96 (suite)
et ceux de l'année précédente ainsi que les principaux habitants. Tous reconnurent que la demande des Duligier était fondée.

Ces détails nous montrent combien le mode d'imposition pouvait fournir aux habitants d'une même localité des occasions de conflit ou des sources d'inimitié. Seuls les exempts pouvaient y échapper, encore les consuls avaient-ils soin de protester à chaque confection de rôle contre de tels privilèges, prétextant que les exempts possédaient précisément les plus beaux fonds de la paroisse. En 1788, ces derniers étaient au nombre de sept : Mr de la Barmondière, Mr Peyson de Bacot, Mr Guillin d'Avenas, Mr de la Roche de Grosbois, Mr de Nully de Lacarelle, Mr le Curé d'Ouroux, Mr Claude Marie Duligier Testenoire curé de la Chapelle sous Dun pour biens patrimoniaux.

Restaient enfin les servis dus à la Seigneurie d'Aloignet. En 1520, au moment où Philibert de Beaujeu et Catherine d'Amboise faisaient dresser leur terrier par le notaire De Fonte, la plupart des fermes appartenaient à des propriétaires particuliers qui vendirent leurs biens aux possesseurs des fiefs de la paroisse. Nous avons vu plus haut quelles étaient, au milieu du XVIIème siècle, les redevances des seigneurs.

Le 9 Décembre 1675, le sieur Jacques Landrin, bourgeois de Lyon, fermier de la Souveraineté des Dombes, Baronnie de Beaujolais et Comté de Châtillon appartenant à son Altesse Royale, Mademoiselle de Montpensier lequel "a soubz affermé comme il soubz afferme à Mr Jean Hugues Testenoire, notaire royal à Ouroux, présent et acceptant, à scavoir la rente noble et seigneurie et dépendances qu'il a dit bien scavoir et ainsy que Claude Morin précédent fermier en a jouy ou dubt jouir et ce pour le temps et terme de six années entières et consécutives commençant au premier janvier prochain, finissant le dernier Décembre 1681 moyennant le prix et somme de 110 livres pour chascune année et la somme de 44 livres pour pot de vin et drouille pour une fois seulement ;

Faict à Villefranche l'an 1675 le 9 Décembre en l'étude de Mr Poyet, notaire. "

% --- Page imprimée 94 — vue 97
De 1696 à 1702, Benoit Morin était fermier de la même rente.
Il prélevait sous le titre de servis la somme de 186 Livres 10 sols
Il ne s'agit pas ici des servis dus à la rente noble de Nagu, ni à celle d'Arcis dont nous n'avons pu retrouver le détail.

Enfin en 1682 Mademoiselle recevait pour "fête balladoire droits de l'aide, jeu de hasard la somme de 15 livres et 2" 4' 6d pour droit de péage.

Par suite du morcellement des propriétés et des transformations qu'elles subirent dans le cours des siècles, les servis dus à la Seigneurie d'Aloignet devinrent d'un calcul fort compliqué. Qu'il nous suffise de citer, vers la fin du siècle dernier les 79/96 parties d'un coupon de seigle et la 48ème partie d'une poule due par la famille Duligier.

Avant de terminer ce chapitre sur la situation de notre paroisse, nous donnons la traduction en langue populaire du bourg d'Ouroux, d'un conte paru dans la revue des Patois N° 2 et qui a pour auteur Mr Combier de Germolles. Il a le mérite d'offrir certaines expressions particulières à nos montagnes et de donner une idée assez exacte du langage usuel. Il peut être intitulé :

Le loup dupé par le renard

Y avait on cou on loup c'avait bien fan. U rincontre on reno que se promenait le lon de la revire. U le deminde c'min i fôdrait fêre pe trôvô à mindzi.

Le reno lé répon :
"Avise, itié, din la revire, cmin y-z-y a de paisson, te nô qu'à mitre ta qué din l'eye, le paisons on fré, i van tüi veni se catsi dedin. Apré, te l'aratzerô é ne von in mindzi n'ton su"

Le lou fit c'min le reno l'avait dé ; é u bot de dieu-z-oure, quin le renô vit qui était bien dzélo, u deci u lou :

Il y avait une fois un loup qui avait bien faim. Il rencontre un renard qui se promenait le long de la rivière. Il lui demande comment il faudrait faire pour trouver à manger.

Le renard lui répond :
"Regarde là dans la rivière, comme il y a de poissons ; tu n'as qu'à mettre ta queue dans l'eau, les poissons ont froid, ils vont tous venir se cacher dedans. Après tu l'arracheras et nous allons en manger notre saoul"

Le loup fit comme le renard l'avait dit et au bout de deux heures, quand le renard vit que c'était bien gelé, il dit au loup :

% --- Page imprimée 95 — vue 98
"Tire, lou, tire; t'in amène de biaux"
Le lou téri tin que sa qué aratzi.
On momin apré y vayent de bredzi in tron de bleuir de tzinda.
"Se ne-z-alin le trovô, que le renô deci u lou, ne les deminderin on batron pe remetre à la plaice de ta qué?
"Té bin encô vré, ce que te di alins y"
Quin le batron fi étatzi, c'min y fayait bien fré, lé bredzi sôtin à la bourde.
Se ne z-y sôtin éri, que le renô deci.
"Ma foi, dze vû bin!
Mé v'la que le féé se méti u batron du lou. Al u biô côri, i ne le tui pô; u contrère. Mé ü courait, mé y brûlait. Se bien q' y le buchlyi c'min on cotzon qu' on vin de tuô.
Le bredzi se metront à creyi: Lô! le lou qu'étô brulo! I ne le tzoume pleu que dou pa su le nô, encore i le fan bien mô!

"Tire, loup, tire, tu en amènes de beaux !"
Le loup tira tant que sa queue fut arrachée.
Un moment après ils voient des bergers en train de teiller du chanvre.
"Si nous allions les trouver, dit le renard au loup, nous leur demanderions une tresse pour remettre à la place de ta queue?
"C'est bien encore vrai ce que tu dis; allons-y"
Quand la tresse fut attachée, comme il faisait bien froid, les bergers sautaient autour d'un feu.
Si nous sautions aussi, lui dit le renard.
"Ma foi je veux bien"
Mais voilà que le feu se mit à la tresse du loup. Il eut beau courir, le feu ne s'éteignit pas au contraire plus il courait, plus il brûlait. Si bien qu'il fut "buclé" comme un porc que l'on vient de tuer.
Les bergers se mettent à crier: "Lô ! le loup qui a été brûlé ! Il ne lui reste plus que deux poils sur le nez, encore ils lui font bien mal !

\begin{center}\textit{[Illustration : sans légende]}\end{center}
